\section{Data Sources}
\label{sec:data-sources}
Technical recruitment often requires multiple data sources to be provided to evaluate candidates beyond traditional CVs, but these are rarely used in practice. These sources have the potential to provide insights into the extensive professional history, technical competencies, behavioural tendencies as well as real-world coding practices. Effective multi-source screening will require a deeper understanding of the strengths and limitations of these data channels.


\begin{table}[htbp]
    \centering
    \caption[Overview of supplementary recruitment data sources]{An overview of supplementary data sources, their value to technical recruitment, and their inclusion in MERIT.}
    \label{tab:data-sources}
    \small
    \renewcommand{\arraystretch}{1.4}
    \begin{tabularx}{\textwidth}{>{\raggedright\arraybackslash\bfseries}p{2.6cm} >{\raggedright\arraybackslash}p{3.2cm} >{\raggedright\arraybackslash}X >{\centering\arraybackslash}p{1.5cm}}
        \hline
        \textbf{Data Source} & \textbf{Primary Function} & \textbf{Value to Technical Recruitment} & \textbf{Included?} \\ \hline
        Curriculum Vitae (CV) & Self-reported professional summary & Baseline for initial screening. Highly vulnerable to embellishment and self-report bias \cite{henle2019_resume_fraud,dunning2004}. & Yes \\
        GitHub & Version-controlled code repositories & Richest source of verifiable work-sample evidence. Demonstrates tech stack usage and coding style \cite{hackerrank2018skills,hackerrank2024skills}. & Yes \\
        LinkedIn & Professional identity and networking & Expands upon CV details. Used by 72\% of recruiters \cite{jobvite2020}. & Yes \\
        LeetCode \& HackerRank & Algorithmic practice and live assessments & Demonstrates tested technical ability and problem-solving under constraints \cite{coderpad2025techhiring}. & No \\
        Stack Overflow & Software engineering Q\&A exchange & Highlights communication style, community engagement, and domain specialisation \cite{vasilescu2013}. & No \\
        Social Media & General personal networking & Used for background checks to identify ``red flags'' \cite{careerbuilder2017_social}. High risk of bias (Note: excludes LinkedIn). & No \\ \hline
    \end{tabularx}
\end{table}

\subsubsection{Relevance to MERIT}
As shown above, there are a variety of external data sources that can be used to evaluate a candidate's suitability for a technical role
beyond the traditional Curriculum Vitae (CV). Being able to capture and integrate multiple data sources can help provide more factors
to differentiate between candidates, allowing for a more comprehensive view of a candidate's suitability for a role. MERIT aims to 
integrate GitHub as a supplementary data source due to it serving as a strong indicator of a candidate's technical ability. LinkedIn
will also be selected as a strong indicator of a candidate's professional experience and also to account for errors in the CV parsing process.
Other data sources such as LeetCode and HackerRank are likely to not be selected as supplementary data sources due to the time 
constraints during development, however, they are still considered as potential future data sources for MERIT.


